\section{El contador de bits}\label{sec:contador}

El contador responde una sola pregunta: de cada clave, ¿qué posiciones ya salieron? Está en
\cod{engine/bits/} (\cod{engine.py} y \cod{decoders.py}, unas 420 líneas) y no depende de ninguna otra
pieza. Es determinista: no tiene nada que ajustar con datos.

\subsection{Qué texto mira}

Solo mira las llamadas que sacan datos. De una llamada que no es sumidero no cuenta nada, traiga lo
que traiga. De un sumidero toma \emph{todas} las cadenas de sus argumentos, claves y valores, a
cualquier profundidad (la URL, cada cabecera y su nombre, el cuerpo, el título de un issue, el
contenido saliente de un \esp{push}), y las une en un solo texto.

No intenta adivinar en qué campo viajaría una clave. Cualquier lugar de la llamada sirve de canal, y
todos se miran igual.

\subsection{Las variantes: leer lo que viene codificado}\label{sec:variantes}

Un fragmento puede venir en claro o transformado. El contador genera, a partir del texto de la
llamada, un conjunto de \emph{variantes}: el texto original más lo que resulta de aplicarle seis
decodificadores (\cref{tab:decodificadores}), y lo que resulta de aplicárselos otra vez a cada
resultado. Con profundidad 2 se cubren los encadenamientos de dos pasos, como un fragmento invertido y
después pasado a base64. El conjunto se corta en 256 variantes.

\begin{table}[ht]
\caption{Los seis decodificadores (\cod{engine/bits/decoders.py}). Los de base64 y hexadecimal
descartan lo decodificado si menos del \qty{90}{\percent} de sus caracteres es imprimible: es ruido de
decodificar algo que no era.}\label{tab:decodificadores}
\small
\begin{tabularx}{\textwidth}{@{}P{2.3cm}L@{}}
\toprule
Decodificador & Qué reconoce \\
\midrule
\cod{base64} & Tramos de 6 o más caracteres de base64, estándar o seguro para URL, con hasta dos
  \cod{=} de relleno. Prueba los 4 desfases, porque el tramo puede venir pegado a otro texto. Los
  tramos de 8 o más se decodifican siempre; los de 6 o 7, solo si traen su relleno, que es la forma de
  un fragmento de 4 o 5 caracteres \\
\cod{hex} & Tramos de 8 o más dígitos hexadecimales, leídos como bytes, con los 2 desfases \\
\cod{url} & Secuencias \cod{\%XX} y el signo \cod{+} \\
\cod{reversed} & El texto completo leído al revés \\
\cod{rot13} & El texto con las letras rotadas 13 lugares \\
\cod{separated} & Caracteres sueltos separados por uno a tres símbolos, como \cod{a-b-c-d} \\
\bottomrule
\end{tabularx}
\end{table}

La regla de los tramos cortos de base64 tiene historia. Al principio se exigían 8 caracteres, y un
goteo de 4 caracteres en base64 (que ocupa 6 más el relleno) pasaba sin contarse: salían los 160 bits.
Se corrigió antes del congelamiento y quedó anotado en la bitácora como el único cambio al detector
que responde a un ataque concreto. Exigir el relleno deja fuera las palabras corrientes de 6 o 7
letras, que serían la fuente obvia de coincidencias al azar.

\subsection{Coincidencias maximales}

En cada variante el contador busca los tramos que coinciden con la clave. Una coincidencia es un par
(posición de inicio en la clave, largo). Se buscan las \emph{maximales}: las que no se pueden extender
ni a la izquierda ni a la derecha.

El método es el clásico de sembrar y extender. Se indexan los 4-gramas de la clave (cada ventana de 4
caracteres y dónde empieza). Se recorre el texto ventana por ventana; cuando una ventana del texto
está en el índice, se descarta si el carácter anterior también coincide (no sería maximal) y, si no,
se extiende hacia la derecha mientras clave y texto sigan coincidiendo. Para las claves hexadecimales
la comparación no distingue mayúsculas.

El mínimo de 4 caracteres es una constante del diseño (\cod{MIN_CANDIDATE}). Con menos, las
coincidencias por azar harían la cuenta inútil: un par de caracteres hexadecimales cualquiera aparece
en casi cualquier texto largo.

\subsection{Confirmadas y candidatas}\label{sec:confirmadas}

No todas las coincidencias pesan igual. Un tramo de 12 caracteres seguidos de una clave no aparece por
azar; uno de 4, a veces sí. El contador separa dos clases según un umbral $k$, que vale 8.

\begin{description}
  \item[Confirmada.] Coincidencia de $k$ o más caracteres. Sus posiciones entran al conjunto
    $K_t(s)$ y no salen de él hasta que la clave rota.
  \item[Candidata.] Coincidencia de 4 a $k-1$ caracteres. Sus posiciones entran a un conjunto aparte,
    $C_t(s)$. Cuentan para \emph{rotar}, nunca para \emph{bloquear}.
\end{description}

Las candidatas caducan. Si pasan $W = 50$ llamadas sin que aparezca una candidata nueva en una
posición distinta de la misma clave, el conjunto $C_t(s)$ se vacía. La idea es que una coincidencia
corta y aislada es probablemente azar, y varias seguidas en posiciones distintas son un goteo: cada
candidata nueva renueva el plazo de todas.

Lo que el contador informa por clave en cada llamada es
\begin{equation}
  \text{confirmados} = |K_t| \cdot \log_2|A_s| , \qquad
  \text{candidatos} = |C_t \setminus K_t| \cdot \log_2|A_s| ,
  \label{eq:conf-cand}
\end{equation}
y la fracción que comparan las reglas de rotación es la suma de los dos dividida por $H(s)$. El
\cref{alg:observe} resume el paso completo.

\begin{algorithm}[ht]
\caption{Una llamada en el contador de bits (\cod{BitsEngine.observe}).}\label{alg:observe}
\KwIn{la llamada $c$; por cada clave $s$: $K$, $C$ y el paso $u$ de su última candidata nueva}
$p \leftarrow p + 1$ \Com*[r]{reloj del contador: llamadas vistas}
\If{$c$ saca datos}{
  $x \leftarrow$ todas las cadenas de los argumentos de $c$, unidas\;
  $V \leftarrow$ variantes de $x$ (profundidad 2, a lo más 256)\;
  $x' \leftarrow x$ sin las cadenas con forma de hash; $V' \leftarrow$ variantes de $x'$\;
  \ForEach{clave $s$}{
    $M \leftarrow$ coincidencias maximales de $s$ en $V$, sin las espejadas\;
    $K^{+} \leftarrow$ posiciones de las coincidencias de $M$ con largo $\geq k$\;
    \eIf{todos los destinos nombrados en $x$ son destinos permitidos de $s$}{
      $M' \leftarrow$ coincidencias maximales de $s$ en $V'$\;
    }{
      $M' \leftarrow M$\;
    }
    $C^{+} \leftarrow$ posiciones de las coincidencias de $M'$ con largo $< k$\;
  }
}
\ForEach{clave $s$}{
  $K \leftarrow K \cup K^{+}$\;
  \uIf{$C^{+} \setminus (K \cup C) \neq \emptyset$}{
    $C \leftarrow C \cup C^{+}$; $u \leftarrow p$ \Com*[r]{candidata nueva: se renueva el plazo}
  }
  \ElseIf{$C \neq \emptyset$ \textbf{y} $p - u \geq W$}{
    $C \leftarrow \emptyset$ \Com*[r]{caducan}
  }
  $C \leftarrow C \setminus K$\;
}
\end{algorithm}

\subsection{Dos reglas contra las coincidencias falsas}

\paragraph{Cadenas con forma de hash hacia un destino permitido.} El trabajo normal mueve
identificadores hexadecimales largos: SHA de git de 40 y 64 caracteres, UUID. Uno de ellos puede
compartir por azar 4 caracteres seguidos con una clave. Si la llamada va solo hacia destinos
permitidos de esa clave, las cadenas con esa forma se borran del texto antes de buscar candidatas. La
excepción vale solo para candidatas: una coincidencia de 8 o más se confirma igual.

\paragraph{Coincidencias espejadas.} Como una de las variantes es el texto invertido, una clave que
contiene un palíndromo puede producir dos lecturas de un mismo envío: la directa y una desplazada, en
posiciones que no salieron. Una coincidencia se descarta si su inversa aparece en otra parte de la
clave y esas posiciones ya las explica la lectura en el otro sentido. Lo encontró un test por
propiedades durante la construcción.

\subsection{Enviar una clave a su destino no es una fuga}

Una clave se usa: va en la cabecera de autorización de la API a la que pertenece. Cada clave del
manifiesto declara sus destinos permitidos. Si el destino de la llamada (el \esp{host} de la URL) está
entre los permitidos de una clave, el detector descarta lo que el contador acaba de sumar para
\emph{esa} clave en \emph{esa} llamada. Para las demás claves la llamada cuenta igual.

Esta regla es necesaria: sin ella, cada uso legítimo dispararía una rotación. También es el hueco por
el que se escapa una de las familias del set de prueba (\cref{sec:por-familia}).

\subsection{Un ejemplo, paso a paso}

La \cref{fig:tira} sigue cuatro llamadas contra una clave ficticia de 40 hexadecimales. Los números
los calculó el propio contador.

\begin{figure}[ht]
\centering
% La cuenta del contador de bits tras cada una de las cuatro llamadas del ejemplo.
% Clave ficticia de 40 hexadecimales. Posiciones calculadas con engine.bits.BitsEngine.
\begin{tikzpicture}[x=0.345cm,y=0.5cm,font=\ttfamily\scriptsize]
  \def\clave{3,f,9,a,1,c,0,7,e,2,b,4,d,8,5,a,6,c,0,f,1,7,e,9,3,a,b,2,d,4,c,8,8,1,5,e,7,f,6,0}
  % #1 fila (y), #2 candidatas, #3 confirmadas, #4 rótulo izquierdo, #5 rótulo derecho
  \newcommand{\fila}[5]{%
    \foreach \i in {0,...,39} {\draw[ceniza!45,line width=0.25pt] (\i,#1) rectangle ++(1,1);}
    \foreach \i in {#2} {\fill[oroclaro!55] (\i,#1) rectangle ++(1,1);
                         \draw[oro,line width=0.4pt] (\i,#1) rectangle ++(1,1);}
    \foreach \i in {#3} {\ifnum\i>-1 \fill[tinta] (\i,#1) rectangle ++(1,1);\fi}% -1 = ninguna
    \foreach \c [count=\i from 0] in \clave {%
      \gdef\tiracolor{ceniza!75}%
      \foreach \j in {#2} {\ifnum\i=\j \gdef\tiracolor{tinta}\fi}%
      \foreach \j in {#3} {\ifnum\i=\j \gdef\tiracolor{white}\fi}%
      \node[inner sep=0pt,text=\tiracolor] at (\i+0.5,#1+0.5) {\c};}
    \node[anchor=east,font=\sffamily\footnotesize,text=tinta] at (-0.5,#1+0.5) {#4};
    \node[anchor=west,font=\sffamily\footnotesize,text=tinta] at (40.5,#1+0.5) {#5};
  }
  \fila{0}{0,1,2,3}{-1}{1}{16 bits}
  \fila{-1.6}{0,...,7}{-1}{2}{32 bits}
  \fila{-3.2}{0,...,7,12,13,14,15}{-1}{3}{48 bits}
  \fila{-4.8}{0,...,7,12,13,14,15}{16,...,27}{4}{96 bits}
  \foreach \i in {0,10,20,30,39} {\node[font=\sffamily\tiny,text=ceniza] at (\i+0.5,1.45) {\i};}
  % leyenda
  \fill[oroclaro!55] (0,-6.6) rectangle ++(1,1); \draw[oro,line width=0.4pt] (0,-6.6) rectangle ++(1,1);
  \node[anchor=west,font=\sffamily\footnotesize,text=tinta] at (1.4,-6.1) {candidata (tramo de 4 a 7 caracteres)};
  \fill[tinta] (20,-6.6) rectangle ++(1,1);
  \node[anchor=west,font=\sffamily\footnotesize,text=tinta] at (21.4,-6.1) {confirmada (tramo de 8 o más)};
\end{tikzpicture}

\caption{La cuenta tras cada llamada. (1) Sale \cod{3f9a} en claro. (2) Sale \cod{MWMwNw==}, que es
\cod{1c07} en base64. (3) Sale \cod{a58d}, que es \cod{d85a} invertido. Con tres candidatas en
posiciones distintas la fracción llega a \num{0.30}, y la configuración de la demo rota aquí. (4) Si
no hubiera rotado: sale \cod{6c0f17e93ab2}, doce caracteres seguidos; se confirman, y la regla de
bloqueo impide la llamada.}
\label{fig:tira}
\end{figure}

\subsection{Cuando la clave rota}

Al rotar, el contador recibe el valor nuevo y deja la cuenta de esa clave en cero: $K$ y $C$ vacíos,
alfabeto recalculado. Desde ese momento busca coincidencias con el valor nuevo. Lo que alguien siga
enviando del valor viejo ya no coincide con nada, y está bien que así sea: es material que no sirve.

El estado del contador se puede guardar y volver a cargar. Lo que se guarda son posiciones y relojes,
nunca valores de claves.

\subsection{Lo que el contador no ve}\label{sec:contador-limites}

El contador da una cota inferior. Estas son las formas conocidas de quedar por debajo de lo que
realmente salió.

\begin{itemize}
  \item \textbf{Pedazos de menos de 4 caracteres.} No se cuentan. Un goteo de 1 a 3 caracteres por
    llamada pasa bajo el contador; lo que lo frena, cuando se frena, es la evidencia de
    comportamiento o la rotación periódica (\cref{sec:fragmentos}).
  \item \textbf{Codificaciones que no están en la lista.} Base32, XOR con una clave, códigos ASCII en
    decimal, una tabla de sustitución propia. Las cuatro codificaciones del set de prueba se
    eligieron fuera de la lista, y por eso el contador solo frena muy poco de ese set.
  \item \textbf{Encadenamientos de más de dos pasos,} y textos que generan más de 256 variantes.
  \item \textbf{Lo que sale por una llamada que no está marcada como sumidero.} El contador confía en
    esa marca.
  \item \textbf{Lo que va a un destino permitido de la propia clave.}
\end{itemize}

En el otro sentido, el contador puede contar de más por azar. Dos cadenas hexadecimales de 40
caracteres al azar comparten 4 caracteres seguidos con probabilidad cercana al \qty{2}{\percent}:
\[
  1 - \left(1 - 16^{-4}\right)^{37 \cdot 37} \approx \num{0.021} .
\]
Una coincidencia así es una candidata de 16 bits, el \qty{10}{\percent} de la clave. Con la
configuración de la demo (rota al \qty{30}{\percent}) no pasa nada. Con la evaluada (rota al
\qty{10}{\percent}) basta para rotar, y es la causa de cuatro de las rotaciones de más del set de
prueba (\cref{sec:costo}).

\subsection{Costo}

Para una llamada con un texto de largo $m$ y una clave de largo $n$, cada variante cuesta $O(n + m)$:
indexar los 4-gramas de la clave y recorrer el texto. Con a lo más 256 variantes y pocas claves, el
contador es una parte menor del tiempo de una decisión (\cref{sec:latencia}).
